\section{Weighted odd-step paths, exact completions, and finite residue actions}
\label{sec:chatnotes-weighted-path-algebra}

This section reconstructs the next coherent programme in the raw
\emph{$Z_n$ Symmetries} export, physical lines 6221--9596.  The controlling
local witness is the chronological export itself.  Two later PDFs entitled
\emph{The Odd-Step Collatz Path Algebra} are recorded as AI-assisted
restatements of parts of that passage, not as scholarly sources and not as
theorem authority.  Their order in the local lineage is
\[
 \begin{array}{c}
 \text{raw chronological passage}\\[-1mm]
 \downarrow\\[-1mm]
 \text{seven-page AI-assisted note}\\[-1mm]
 \downarrow\\[-1mm]
 \text{twenty-seven-page AI-assisted note}
 \end{array}
\]
The passage contains a substantial finite mathematical core, but it also
uses mutually incompatible composition orders, calls nonunital finite
semigroups monoids, puts infinite sums in a finite-support algebra, and
identifies a prefix $3$-adic cocycle with a forward affine intercept without
the required word reversal.  We retain the objects and calculations at their
proved scope and repair each of those defects explicitly.

The general algebraic vocabulary is crosswalked against the standard
definitions of free monoids, semirings, formal power series, weighted paths,
and transformation monoids in
\cite{KuichSalomaa1986,DrosteKuichVogler2009,Steinberg2016}.  Those sources do
not contain the Collatz-specific theorems below; every such theorem is proved
here.

\subsection{The odd first-return category and its local block category}

Let
\[
 \mathcal O=\{1,3,5,\ldots\},
 \qquad
 a(x)=\nu_2(3x+1),
 \qquad
 C(x)=\frac{3x+1}{2^{a(x)}}
 \quad(x\in\mathcal O).
\]
Thus $C$ is the positive-odd first-return map.  Write
\[
 o(n)=\frac{3n+1}{2}\quad(n\text{ odd}),
 \qquad
 e(n)=\frac n2\quad(n\text{ even})
\]
for the two legal branches of the shortened map, with their domains retained.

\begin{definition}[first-return path category]
\label{def:odd-first-return-path-category}
The category $\mathcal E_C$ has object set $\mathcal O$.  For every
$x\in\mathcal O$ it has a generating arrow
\[
 \epsilon_x:x\longrightarrow C(x),
 \qquad \lambda(\epsilon_x)=a(x),
\]
and its morphisms are finite composable strings of these arrows, including
the empty path $1_x$ at each object.  A path
\[
 p:x_0\xrightarrow{a_1}x_1\xrightarrow{a_2}\cdots
       \xrightarrow{a_n}x_n
\]
has length and total exponent
\[
 \ell(p)=n,
 \qquad A(p)=a_1+\cdots+a_n,
\]
and weight
\[
 w(p)=u^{\ell(p)}v^{A(p)-\ell(p)}.
\]
The category $\mathcal D_C^{\mathrm{blk}}$ has the same objects.  Its
morphisms are precisely the legal shortened-map paths between odd endpoints
whose local word has the form
\[
 (oe^{m_1})(oe^{m_2})\cdots(oe^{m_n}),
 \qquad m_i\geq0.
\]
For a block path $q$, put
\[
 \#_o(q)=\text{number of $o$-steps},
 \qquad
 \#_e(q)=\text{number of $e$-steps},
 \qquad
 w_{\mathrm{blk}}(q)=u^{\#_o(q)}v^{\#_e(q)}.
\]
Composition in both categories is path concatenation.
\end{definition}

The set of strict first-return edges by itself is a directed graph, not a
category: a composite of two strict first-return edges is usually not a
single strict first-return edge.  The free path category $\mathcal E_C$ is the
typed categorical object used below.

\begin{theorem}[exact odd-block category isomorphism]
\label{thm:chatnotes-block-category-isomorphism}
For a positive exponent put
\[
 \iota(a)=oe^{a-1}.
\]
There is a functor
\[
 \iota:\mathcal E_C\longrightarrow\mathcal D_C^{\mathrm{blk}}
\]
which is the identity on objects and replaces each exponent edge by its local
block.  It is an isomorphism.  Its inverse is
\[
 U\bigl((oe^{m_1})\cdots(oe^{m_n})\bigr)
   =(m_1+1,\ldots,m_n+1).
\]
Moreover
\[
 \#_o(\iota(p))=\ell(p),
 \qquad
 \#_e(\iota(p))=A(p)-\ell(p),
 \qquad
 w_{\mathrm{blk}}(\iota(p))=w(p).
\]
In particular, this block passage loses no local trajectory information.
\end{theorem}

\begin{proof}
Fix $x\in\mathcal O$ and put $a=a(x)$.  The first local step gives
$(3x+1)/2$.  If $a>1$, this value and the following $a-2$ intermediate
values are even, so exactly $a-1$ legal $e$-steps follow.  Their endpoint is
$(3x+1)/2^a=C(x)$, which is odd.  Conversely, if a legal local path from an
odd $x$ has word $oe^m$ and first returns to an odd state at its endpoint,
then $3x+1$ is divisible by $2^{m+1}$ and not by $2^{m+2}$.  Hence
$a(x)=m+1$.

Every displayed block word has a unique parsing: each occurrence of $o$
starts a block and every following run of $e$'s ends immediately before the
next $o$ or at the word endpoint.  Therefore $U\iota$ and $\iota U$ are the
identity on every morphism.  Both maps preserve endpoints and concatenation,
so they are mutually inverse functors.  Counting the one $o$ and the $a_i-1$
copies of $e$ in each block proves the three weight identities.
\end{proof}

The theorem corrects a recurrent description of $U$ as a proper forgetful
quotient.  It would be information-losing on an unrestricted local word
category, but on the explicitly defined legal odd-to-odd block category it
has the displayed two-sided inverse.

\subsection{Two different \texorpdfstring{$\mathbb F_2^2$}{F2-squared} degrees}

The weighted block counts give a degree
\[
 d_{\mathrm{count}}(p)
 =\bigl(\ell(p)\bmod2,\,(A(p)-\ell(p))\bmod2\bigr).
\]
The coefficient-label convention in
Section~\ref{sec:chatnotes-coefficient-algebra}, which assigns $00$ to the
even branch and $11$ to the odd branch, instead gives
\[
 d_{\mathrm{coeff}}(p)
   =\bigl(\ell(p)\bmod2,\ell(p)\bmod2\bigr).
\]

\begin{proposition}[exact comparison of the two degrees]
\label{prop:chatnotes-two-path-gradings}
The homomorphism
\[
 \kappa:\mathbb F_2^2\longrightarrow\mathbb F_2^2,
 \qquad \kappa(r,s)=(r,r),
\]
satisfies $d_{\mathrm{coeff}}=\kappa\circ d_{\mathrm{count}}$.  Its kernel,
image, and fibres are
\[
 \ker\kappa=\{00,01\},
 \qquad \operatorname{im}\kappa=\{00,11\},
\]
\[
 \kappa^{-1}(00)=\{00,01\},
 \qquad
 \kappa^{-1}(11)=\{10,11\},
\]
with empty fibres over $01$ and $10$.  Thus $\kappa$ is an idempotent
projection onto the diagonal $\Delta=\{00,11\}$.  Equivalently, its codomain
restriction induces the quotient isomorphism
$\mathbb F_2^2/\ker\kappa\simeq\Delta$.  It is not a quotient onto the whole
displayed codomain, and the two degrees are not interchangeable.
\end{proposition}

\begin{proof}
Substitution gives
$\kappa(\ell,A-\ell)=(\ell,\ell)$ modulo two.  The kernel, image, and four
fibres follow by evaluating the four elements of $\mathbb F_2^2$.
\end{proof}

\subsection{Composition order and the ordinary category algebra}

We use the standard target-left convention throughout this section.  If
$p:x\to y$ and $q:y\to z$, then $q\circ p$ means first $p$, then $q$, and
the category-algebra product is
\begin{equation}
\label{eq:weighted-path-target-left-product}
 [q][p]=[q\circ p],
 \qquad
 [q][p]=0\quad\text{when }t(p)\ne s(q).
\end{equation}
Consequently
\[
 [p]=e_{t(p)}[p]e_{s(p)},
 \qquad e_x=[1_x],
\]
and for the path in Definition~\ref{def:odd-first-return-path-category},
\[
 [p]=[\epsilon_{x_{n-1}}]\cdots[\epsilon_{x_0}].
\]
Chronological word concatenation is the opposite written order from this
target-left product.  Several raw formulas use one order in the path algebra
and the other order in a transformation or pullback algebra without an
opposite construction; Subsection~\ref{subsec:typed-finite-actions} gives the
forced repair.

Let $R$ be a commutative ring.  The ordinary weighted category algebra is
\[
 \mathfrak A_C
   =\bigoplus_{p\in\operatorname{Mor}(\mathcal E_C)}R[u,v][p]
\]
with multiplication \eqref{eq:weighted-path-target-left-product}.  Every
element has finite path support.  Therefore neither the sum of all object
idempotents nor the sum of all finite paths is generally an element of
$\mathfrak A_C$.

\subsection{Endpoint series and the exact path completion}

For $x,y\in\mathcal O$, define
\begin{equation}
\label{eq:weighted-endpoint-series}
 Z_{x\to y}(u,v)
   =\sum_{\substack{n\geq0\\ C^n(x)=y}}
      u^n v^{A_n(x)-n}
   \in R[v][[u]],
\end{equation}
where $A_0(x)=0$ and
$A_n(x)=\sum_{j=0}^{n-1}a(C^j(x))$.  The coefficient of each $u^n$ is
either zero or one monomial in $v$; no convergence assertion is needed for
this formal series.

\begin{proposition}[exact hit-time alternatives]
\label{prop:weighted-endpoint-hit-times}
For fixed $x,y\in\mathcal O$, the set
\[
 H_{x,y}=\{n\geq0:C^n(x)=y\}
\]
is empty, a singleton, or
\[
 r+d\mathbb N_0
\]
for a unique first hit $r$ and the least positive period $d$ of $y$.  In the
last case, if
\[
 B=\sum_{j=0}^{d-1}a(C^j(y)),
\]
then
\begin{equation}
\label{eq:weighted-cyclic-endpoint-rational-form}
 Z_{x\to y}(u,v)
 =\frac{u^r v^{A_r(x)-r}}{1-u^d v^{B-d}}
\end{equation}
as a formal geometric series.  Every positive $d$-cycle satisfies
\begin{equation}
\label{eq:positive-cycle-weight-inequality}
 2^B>3^d.
\end{equation}
Hence the numerical specialization $(u,v)=(3/2,1/2)$ converges on every
cyclic endpoint series.  In contrast,
\[
 Z_{1\to1}(1,1)=\sum_{n\geq0}1
\]
diverges.
\end{proposition}

\begin{proof}
If $m<n$ are two hit times, then $C^{n-m}(y)=y$, so $y$ is periodic.  Once
the orbit first reaches a periodic $y$, its later returns occur exactly at
multiples of the least period.  This proves the alternatives and the
uniqueness of $r,d$.  Each traversal of the cycle adds $d$ to the length and
$B$ to the exponent sum, which gives
\eqref{eq:weighted-cyclic-endpoint-rational-form}.

Write the cycle as $y_i\xrightarrow{a_i}y_{i+1}$ with indices modulo $d$.
For every $i$,
\[
 2^{a_i}y_{i+1}=3y_i+1>3y_i.
\]
Multiplication and cancellation of the positive product $\prod_i y_i$
give $2^B>3^d$.  At $(3/2,1/2)$ the geometric ratio is $3^d/2^B<1$.
Finally $C(1)=1$ and $a(1)=2$, so every term of
$Z_{1\to1}$ has weight $(uv)^n$.
\end{proof}

The completion needed for global path sums is not the preceding direct sum.

\begin{definition}[coefficientwise path completion]
\label{def:weighted-path-product-completion}
Put
\[
 \widehat{\mathfrak A}_C
  =\prod_{p\in\operatorname{Mor}(\mathcal E_C)}R[u,v][p].
\]
For $F=(F_p)$ and $G=(G_p)$, define the coefficient at a path $r$ by
\[
 (FG)_r=\sum_{r=q\circ p}F_qG_p.
\]
This is a finite sum: a path of length $n$ has exactly $n+1$ cuts.
\end{definition}

The finite-cut statement proves that multiplication is well defined for
arbitrary path support.  Associativity follows by identifying both
parenthesizations with the finite set of two-cut decompositions of a fixed
path.  The completion contains
\[
 E=\sum_{x\in\mathcal O}e_x,
 \qquad
 W=\sum_{x\in\mathcal O}u v^{a(x)-1}[\epsilon_x],
 \qquad
 G=\sum_{p}w(p)[p].
\]
Every nonempty path has a unique first edge and a unique last edge, so
\begin{equation}
\label{eq:weighted-path-resolvent-identities}
 G=E+WG=E+GW.
\end{equation}
Because $R$ is a ring, this says $G=(E-W)^{-1}$ inside the displayed
completion.  If $S$ is instead a commutative semiring, repeat
Definition~\ref{def:weighted-path-product-completion} with coefficients
$S[u,v]$.  The finite-cut convolution and the identities
$G=E+WG=E+GW$ remain valid there and define the Kleene star $W^*$; the
subtraction and inverse notation are not used.  None of $E,W,G$ is silently
placed in an ordinary finite-support algebra.

\subsection{The predecessor transfer operator}

For odd $y$ define
\[
 \mathcal P(y)=\left\{a\geq1:
       x_a(y)=\frac{2^a y-1}{3}\text{ is an integer}\right\}.
\]

\begin{lemma}[exact predecessor fibres]
\label{lem:weighted-predecessor-fibres}
For odd $y$,
\[
 \mathcal P(y)=
 \begin{cases}
  \varnothing,&y\equiv0\pmod3,\\
  \{2,4,6,\ldots\},&y\equiv1\pmod3,\\
  \{1,3,5,\ldots\},&y\equiv2\pmod3.
 \end{cases}
\]
For every $a\in\mathcal P(y)$, $x_a(y)$ is positive odd and
\[
 C(x_a(y))=y,
 \qquad a(x_a(y))=a.
\]
Every positive odd predecessor of $y$ occurs exactly once in this list.
\end{lemma}

\begin{proof}
Modulo three, $2^a y=1$ is impossible when $3\mid y$.  Since $2^a$ is
$1$ for even $a$ and $2$ for odd $a$ modulo three, the other two cases have
the stated parities.  The numerator is positive and odd; division by three
preserves oddness.  The identity
$3x_a(y)+1=2^ay$ with $y$ odd proves both displayed first-return identities.
Solving that same identity for $x$ proves uniqueness.
\end{proof}

Define
\begin{equation}
\label{eq:weighted-predecessor-transfer}
 (Lf)(y)=u\sum_{a\in\mathcal P(y)}v^{a-1}f(x_a(y)).
\end{equation}
This formula has several distinct exact domains.

\begin{theorem}[algebraic, formal, and analytic transfer domains]
\label{thm:weighted-transfer-exact-domains}
Let $R$ be a commutative ring.

\begin{enumerate}[label=\textup{(\roman*)}]
\item Formula \eqref{eq:weighted-predecessor-transfer} defines an
$R[u,v]$-linear endomorphism of the finite-support module
$R[u,v]^{(\mathcal O)}$.

\item It defines a continuous endomorphism of
\[
 \mathcal M=(R[[v]][[u]])^{\mathcal O}.
\]
All $v$-series in this product are interpreted coefficientwise.  The module
$\mathcal M$ carries the complete linear, uniform $u$-adic topology with
\[
 F^k\mathcal M=(u^kR[[v]][[u]])^{\mathcal O}\qquad(k\geq0).
\]
For a fixed coefficient of $v^j$, only $a\leq j+1$ can contribute, and $L$
obeys $L(F^k\mathcal M)\subseteq F^{k+1}\mathcal M$.  Consequently
\[
 (I-L)^{-1}=\sum_{n\geq0}L^n
\]
is well defined $u$-adically and
\begin{equation}
\label{eq:transfer-endpoint-kernel}
 \bigl[(I-L)^{-1}\delta_x\bigr](y)=Z_{x\to y}(u,v).
\end{equation}

\item Let $u,v\in\mathbb C$ and $u\ne0$.  On counting-measure sequence
spaces, $L$ is bounded on $\ell^1(\mathcal O)$ exactly when $|v|\leq1$, and
then
\[
 \lVert L\rVert_{\ell^1\to\ell^1}=|u|.
\]
For $1<p\leq\infty$, let $q$ be conjugate to $p$, with $q=1$ when
$p=\infty$.  Then $L$ is bounded on $\ell^p(\mathcal O)$ exactly when
$|v|<1$, and
\begin{equation}
\label{eq:weighted-transfer-lp-norm}
 \lVert L\rVert_{\ell^p\to\ell^p}
 =|u|\bigl(1-|v|^{2q}\bigr)^{-1/q}.
\end{equation}
The same strict condition and norm hold on $c_0(\mathcal O)$.
\end{enumerate}
If $u=0$, the zero operator is defined for every finite $v$.
\end{theorem}

\begin{proof}
For (i), a finitely supported $f$ has only finitely many states $x$ with
$f(x)\ne0$, and each such state has one image $C(x)$.  Thus the apparently
infinite predecessor formula is a finite sum on this domain.

For (ii), the coefficient of $v^j$ can receive a term only from $a-1\leq j$.
This proves coefficientwise finiteness.  The factor $u$ proves the filtration
statement and hence the Neumann identity.  Starting from $\delta_x$, each
application of $L$ follows the unique forward edge from the current odd
state and multiplies by $uv^{a-1}$.  Induction gives
\[
 (L^n\delta_x)(y)
 =\begin{cases}
 u^n v^{A_n(x)-n},&C^n(x)=y,\\
 0,&C^n(x)\ne y,
 \end{cases}
\]
which proves \eqref{eq:transfer-endpoint-kernel}.

For (iii), the input coordinates split into disjoint fibres $C^{-1}(y)$.
The absolute row weights over a fibre are either
\[
 |u|(1,|v|^2,|v|^4,\ldots)
 \quad\text{or}\quad
 |u|(|v|,|v|^3,|v|^5,\ldots),
\]
or the empty sequence.  For $1<p\leq\infty$, Hölder's inequality gives the
row-functional norm as the $\ell^q$ norm of this weight sequence.  The first
row is the larger one, and its norm is exactly the right side of
\eqref{eq:weighted-transfer-lp-norm}.  Finite truncations with phases aligned
to the weights approach equality, so the upper bound is sharp.  If
$|v|\geq1$, those same finite truncations have unbounded row-functional norm.

For $p=1$, each input state occupies exactly one column and has column sum
$|u||v|^{a(x)-1}$.  All positive exponents occur by
Lemma~\ref{lem:weighted-predecessor-fibres}.  The supremum is $|u|$ when
$|v|\leq1$ and is infinite when $|v|>1$.  Finally, for $|v|<1$, the same row
estimate maps $c_0$ to itself: split the series into finitely many exponents,
whose predecessors tend to infinity with $y$, and a uniformly small
geometric tail.  Finite-support phase-aligned vectors prove sharpness.  When
$|v|\geq1$, choose values tending slowly to zero along one predecessor fibre
and align their phases; the defining row sum fails to be bounded.  This proves
the $c_0$ assertion.
\end{proof}

At $u=v=1$, the operator is therefore norm one on $\ell^1$, but it is
unbounded on $\ell^\infty$ and is not defined on the constant function one by
the raw infinite-sum formula.  An arbitrary product of polynomial-valued
functions is likewise not an algebraic domain for $L$: the missing $v$-adic
completion is essential.

\subsection{Prefix cocycles and forward affine intercepts}

For a finite exponent word $w=a_1\cdots a_n$, put
\[
 A_k(w)=a_1+\cdots+a_k,
 \qquad
 A(w)=A_n(w),
\]
and define the finite prefix coordinate
\begin{equation}
\label{eq:weighted-prefix-cocycle}
 \rho(w)=\sum_{k=1}^n3^{k-1}2^{-A_k(w)}\in\mathbb Z_3.
\end{equation}
For the empty word, $A(\varnothing)=0$ and $\rho(\varnothing)=0$.

\begin{proposition}[typed prefix cocycle and its exact fibres]
\label{prop:weighted-prefix-cocycle-injection}
For exponent words $w,z$ in chronological concatenation order,
\begin{equation}
\label{eq:weighted-prefix-concatenation}
 \rho(wz)=\rho(w)+3^{|w|}2^{-A(w)}\rho(z).
\end{equation}
Let
\[
 P=\{(m,A)\in\mathbb N_0^2:A\geq m\},
 \qquad
 \theta_{(m,A)}(r)=3^m2^{-A}r\quad(r\in\mathbb Z_3).
\]
With addition on $P$ and multiplication
\[
 (r,p)\star(s,q)=(r+\theta_p(s),p+q),
\]
the semidirect product $\mathbb Z_3\rtimes_\theta P$ is a monoid and
\[
 \chi_{\mathrm{pre}}(w)=(\rho(w),(|w|,A(w)))
\]
is an injective monoid homomorphism from positive-exponent words.

More explicitly, put $R(w)=2^{A(w)}\rho(w)\in\mathbb Z$.  If $n>1$, then
\begin{equation}
\label{eq:weighted-prefix-recursive-inverse}
 a_n=\nu_2\bigl(R(w)-3^{n-1}\bigr),
 \qquad
 R(a_1\cdots a_{n-1})
 =\frac{R(w)-3^{n-1}}{2^{a_n}}.
\end{equation}
Together with $a_1=A(w)$ when $n=1$, this is a two-sided recursive inverse
on the image.  Every fibre of $\chi_{\mathrm{pre}}$ is therefore a singleton.
\end{proposition}

\begin{proof}
Splitting the sum \eqref{eq:weighted-prefix-cocycle} after $|w|$ terms and
factoring $3^{|w|}2^{-A(w)}$ from the remainder proves
\eqref{eq:weighted-prefix-concatenation}.  The endomorphisms satisfy
$\theta_p\theta_q=\theta_{p+q}$, so the displayed semidirect multiplication is
associative with identity $(0,(0,0))$; the cocycle identity proves
multiplicativity of $\chi_{\mathrm{pre}}$.

Multiplying \eqref{eq:weighted-prefix-cocycle} by $2^A$ gives
\[
 R(w)=\sum_{k=1}^n3^{k-1}2^{A-A_k}.
\]
The last term is $3^{n-1}$.  After subtracting it, the preceding term has
exact $2$-adic valuation $a_n$, while every earlier term is divisible by
$2^{a_{n-1}+a_n}$.  Hence no cancellation changes that valuation, proving
the first formula in \eqref{eq:weighted-prefix-recursive-inverse}.  Division
gives the same integer for the prefix word.  Recursion recovers all exponents,
which proves injectivity and the inverse identities.
\end{proof}

The acting object here is a monoid, not a group: $\theta_{(m,A)}$ is not
surjective on $\mathbb Z_3$ when $m>0$.  In the coordinate pair $(m,A)$ its
multiplier is $3^m2^{-A}$, not $3^m2^{-(m+A)}$.

The preceding injectivity is a statement about exponent words.  Erasing the
source and target of a path is a different operation and has exact infinite
fibres.

\begin{proposition}[exact source fibres of an exponent word]
\label{prop:weighted-word-source-fibres}
For a word $w=(a_1,\ldots,a_n)$ with $n\geq0$, put $A=A(w)$ and
\[
 B_w=\sum_{i=1}^n3^{n-i}2^{a_1+\cdots+a_{i-1}}\in\mathbb Z.
\]
The word $w$ is the exact successive valuation word of a positive odd source
$x$ if and only if
\begin{equation}
\label{eq:weighted-word-source-congruence}
 x\equiv3^{-n}(2^A-B_w)\pmod{2^{A+1}}.
\end{equation}
Thus every exponent word has one odd residue class of positive sources modulo
$2^{A+1}$, hence infinitely many path-category realizations.  Retaining the
source object selects one of them.
\end{proposition}

\begin{proof}
If $n=0$, the claim is the tautology $x\equiv1\pmod2$ for a positive odd
source, and the word is empty.  Now let $n\geq1$.  We prove directly, by
induction on $n$, that $w$ is the exact valuation
word of $x$ if and only if
\begin{equation}
\label{eq:weighted-word-terminal-congruence}
 3^nx+B_w\equiv2^A\pmod{2^{A+1}}.
\end{equation}
For $n=1$, this is
$3x+1\equiv2^{a_1}\pmod{2^{a_1+1}}$, exactly the assertion
$\nu_2(3x+1)=a_1$.

For $n>1$, write $w'=(a_2,\ldots,a_n)$,
$A'=A-a_1$, and let $B_{w'}$ be its displayed integer.  Then
\[
 B_w=3^{n-1}+2^{a_1}B_{w'}.
\]
Reducing \eqref{eq:weighted-word-terminal-congruence} modulo $2^{a_1}$
first shows that $2^{a_1}$ divides $3x+1$.  Put
$x_1=(3x+1)/2^{a_1}$.  Dividing the full congruence by $2^{a_1}$ gives
\[
 3^{n-1}x_1+B_{w'}\equiv2^{A'}\pmod{2^{A'+1}}.
\]
Here $A'\geq1$ and $B_{w'}$ is odd, so reduction modulo two shows that
$x_1$ is odd; positivity is immediate from $x>0$.  By induction the displayed
congruence is equivalent to $x_1$ being positive odd with exact
valuation word $w'$.  In particular $x_1$ is odd, so the first valuation is
exactly $a_1$.  Reversing the argument proves the converse.

Solving \eqref{eq:weighted-word-terminal-congruence} for $x$ is legitimate
because $3^n$ is a unit modulo $2^{A+1}$, and gives the unique class
\eqref{eq:weighted-word-source-congruence}.  For $n\geq1$, the integer $B_w$
is odd: its first summand is odd and every later summand is even.  Hence the
class is odd; the $n=0$ class was handled separately.
Adding positive multiples of $2^{A+1}$ gives infinitely many positive
sources.  This direct induction also supplies, without using it as a proof
dependency, the finite characterization appearing in Tao's source
\cite[source lines 455--470]{Tao2026v7}.
\end{proof}

Now put
\[
 g_a(X)=\frac{3X+1}{2^a}
\]
and apply $a_1$ first.  The forward chronological composite is
\begin{equation}
\label{eq:weighted-forward-affine-word}
 G_w=g_{a_n}\circ\cdots\circ g_{a_1}
 =3^n2^{-A(w)}X+\beta(w),
\end{equation}
where
\begin{equation}
\label{eq:weighted-forward-beta}
 \beta(w)=\sum_{i=1}^n3^{n-i}2^{-(a_i+\cdots+a_n)}.
\end{equation}
Thus $B_w=2^{A(w)}\beta(w)$, now as a consequence of the two displayed
formulas rather than as a definition used before $\beta$ is introduced.

\begin{proposition}[exact reversal to Tao's offset]
\label{prop:weighted-prefix-forward-reversal}
For every finite positive-exponent word,
\[
 \beta(w)=\rho(w^{\mathrm{rev}}).
\]
Moreover $\beta(w)$ is exactly Tao's finite offset $F_n(a_1,\ldots,a_n)$ in
the forward convention of \cite[source lines 208--229 and
384--393]{Tao2026v7}.  Thus the prefix coordinate and forward intercept are
connected by word reversal; they are not equal on the same chronological
word in general.
\end{proposition}

\begin{proof}
Reverse the index in \eqref{eq:weighted-prefix-cocycle}.  The $k$th prefix
sum of $w^{\mathrm{rev}}$ is the $k$-term suffix sum of $w$, and its
coefficient $3^{k-1}$ becomes $3^{n-i}$.  This is exactly
\eqref{eq:weighted-forward-beta}.  Induction in
\eqref{eq:weighted-forward-affine-word} gives the same sum, and direct
comparison with Tao's displayed definition gives the final identity.  This
also agrees with the pointwise reversal already proved in
Proposition~\ref{prop:Siegel-Tao-Syracuse-coordinate}.
\end{proof}

For example,
\[
 \rho(1,4)=\rho(3,1,1)=\frac{19}{32},
 \qquad
 \beta(1,4)=\frac5{32},
 \qquad
 G_{(1,4)}(X)=\frac{9X+5}{32}.
\]
The distinction is not cosmetic: $G_{(1,4)}(3)=1$, whereas the unreversed
prefix constant would give the wrong value.  The first equality also shows
that $\rho$ alone is not injective.  The decoration in $\chi_{\mathrm{pre}}$
is essential: the two words have respective $(|w|,A(w))$ coordinates $(2,5)$
and $(3,5)$.

Because target-left category multiplication writes later arrows on the left,
$\chi_{\mathrm{pre}}$ becomes covariant on the path algebra only after its
coefficient monoid is replaced by the opposite monoid.  This typing is made
explicit in Proposition~\ref{prop:weighted-typed-representations} below.

\subsection{Finite residue transformation semigroups}

Fix $N\geq1$ and write
\[
 R_N=\mathbb Z/3^N\mathbb Z,
 \qquad U_N=R_N^\times,
 \qquad \varphi_N=|U_N|=2\cdot3^{N-1}.
\]
Since
\[
 \operatorname{ord}_{3^N}(2)=2\cdot3^{N-1},
\]
the elements $2^{-a}$, $a\geq1$, run through all of $U_N$.  For
$\eta\in U_N$ put
\[
 G_\eta(r)=\eta(3r+1)=3\eta r+\eta.
\]
Let $S_N$ be the image of all nonempty words in the transformations $G_\eta$,
and let
\[
 T_N=S_N\sqcup\{\operatorname{id}_{R_N}\}.
\]
Every nonempty word sends $0$ to a unit, so the identity is not in $S_N$.
Thus $S_N$ is a transformation semigroup and $T_N$ is its identity-adjoined
transformation monoid.

For chronological $\boldsymbol\eta=(\eta_1,\ldots,\eta_k)$, with $\eta_1$
applied first,
direct expansion gives
\begin{equation}
\label{eq:finite-residue-word-affine-formula}
 G_{\boldsymbol\eta}(r)
 =3^k\left(\prod_{i=1}^k \eta_i\right)r
  +\sum_{j=1}^k3^{k-j}\prod_{i=j}^k \eta_i.
\end{equation}
The constant term is congruent to $\eta_k$ modulo three and is therefore a
unit.

\begin{lemma}[order of two modulo powers of three]
\label{lem:order-two-mod-three-power}
For every $N\geq1$,
$\operatorname{ord}_{3^N}(2)=2\cdot3^{N-1}$.
\end{lemma}

\begin{proof}
By Euler's theorem the order divides
$\varphi(3^N)=2\cdot3^{N-1}$, and reduction modulo three shows that it is
even.  Hence it has the form $2\cdot3^j$ with $0\leq j\leq N-1$.  The
lifting-the-exponent
identity gives
\[
 \nu_3(2^{2\cdot3^j}-1)
 =\nu_3(4^{3^j}-1)=\nu_3(4-1)+\nu_3(3^j)=j+1.
\]
The congruence $2^{2\cdot3^j}\equiv1\pmod{3^N}$ therefore holds exactly when
$j\geq N-1$.  The only possible order is $2\cdot3^{N-1}$; the case $N=1$
is included with $j=0$.
\end{proof}

For $N\geq2$ define
\[
 L_{N,1}=\{r\mapsto3br+b:b\in U_N\},
\]
and for $2\leq k<N$ define
\[
 L_{N,k}=\{r\mapsto3^ksr+b:
          s\in(\mathbb Z/3^{N-k}\mathbb Z)^\times, b\in U_N\}.
\]
Finally let
\[
 \mathcal C_N=\{c_b:b\in U_N\},
 \qquad c_b(r)=b.
\]

\begin{theorem}[exact finite residue normal form]
\label{thm:finite-residue-semigroup-normal-form}
For $N=1$,
\[
 S_1=\mathcal C_1=\{c_1,c_2\}.
\]
For every $N\geq2$, there is a disjoint union
\begin{equation}
\label{eq:finite-residue-semigroup-strata}
 S_N=L_{N,1}\sqcup
     \bigsqcup_{k=2}^{N-1}L_{N,k}\sqcup\mathcal C_N.
\end{equation}
More precisely, $L_{N,k}$ is exactly the set represented by words of length
$k<N$, every word of length at least $N$ is constant, and every unit-valued
constant already has a representative of length $N$.
\end{theorem}

\begin{proof}
Formula \eqref{eq:finite-residue-word-affine-formula} shows that a word of
length $k<N$ has slope of exact $3$-adic valuation $k$ and unit intercept;
at length at least $N$ it is constant.  At length one, slope and intercept
satisfy the exceptional relation $\alpha=3b$, giving $L_{N,1}$.

We prove surjectivity onto every later stratum constructively.  First suppose
$N\geq3$.  At length two, prescribe
$s\in(\mathbb Z/3^{N-2}\mathbb Z)^\times$ and $b\in U_N$.
Put
\[
 q=sb^{-1}\pmod{3^{N-2}},
 \qquad
 \eta_1=\frac q{1-3q}\pmod{3^{N-2}},
\]
choose any unit lift of $\eta_1$ to $U_N$, and set
\[
 \eta_2=\frac b{1+3\eta_1}\pmod{3^N}.
\]
Then
\[
 \eta_2(3\eta_1+1)=b,
 \qquad \eta_2\eta_1\equiv s\pmod{3^{N-2}},
\]
so $G_{\eta_2}\circ G_{\eta_1}(r)=9sr+b$.  When $N=2$, choose an arbitrary
$\eta_1\in U_2$ and put $\eta_2=b/(1+3\eta_1)$ in $U_2$; the resulting
length-two word is the constant $c_b$.

Suppose all length-$(k-1)$ forms occur, where $3\leq k<N$.  Given a desired
$3^ks'r+b'$, choose any prior unit intercept $b_0$, put
\[
 \eta_k=\frac{b'}{3b_0+1},
\]
and choose the prior slope parameter $s_0$ so that
$s_0\equiv s'\eta_k^{-1}\pmod{3^{N-k}}$.  Applying $G_{\eta_k}$ on the left
gives the prescribed slope and intercept.  To obtain a prescribed constant
$c_{b'}$ at length $N\geq3$, choose any already constructed length-$(N-1)$
form with intercept $b_0$ and apply
$G_{\eta_N}$ with $\eta_N=b'/(3b_0+1)$; its slope vanishes modulo $3^N$.
Equality of two affine maps is equivalent to equality of their
values at $0$ and $1$, hence of intercept and slope; distinct strata have
different slope valuations.  This proves both completeness and disjointness.
\end{proof}

\begin{corollary}[cardinalities and ranks]
\label{cor:finite-residue-cardinalities-ranks}
For $N\geq2$,
\begin{equation}
\label{eq:finite-residue-cardinality}
 |S_N|=\varphi_N(3^{N-2}+1),
 \qquad
 |T_N|=1+\varphi_N(3^{N-2}+1).
\end{equation}
The first six semigroup cardinalities are
\[
 2,\ 12,\ 72,\ 540,\ 4536,\ 39852,
\]
and the corresponding monoid cardinalities are
\[
 3,\ 13,\ 73,\ 541,\ 4537,\ 39853.
\]
An element of $L_{N,k}$ has image $b+3^kR_N$ and rank $3^{N-k}$;
the constants have rank one, while the adjoined identity has rank $3^N$.
\end{corollary}

\begin{proof}
The stratum sizes are
\[
 |L_{N,1}|=\varphi_N,
 \quad |L_{N,k}|=\varphi_N\varphi_{N-k}\ (2\leq k<N),
 \quad |\mathcal C_N|=\varphi_N.
\]
Since $\sum_{j=1}^{N-2}\varphi_j=3^{N-2}-1$, their sum is
\eqref{eq:finite-residue-cardinality}.  The image and rank statements follow
by varying $3^ksr$ over the ideal $3^kR_N$.
\end{proof}

\begin{proposition}[the minimum ideal]
\label{prop:finite-residue-minimum-ideal}
The rank-one stratum $\mathcal C_N$ is the minimum two-sided ideal of both
$S_N$ and $T_N$.  Its multiplication is left-zero:
\[
 c_bc_d=c_b.
\]
\end{proposition}

\begin{proof}
For every $f\in T_N$ and unit $b$,
\[
 c_bf=c_b,
 \qquad fc_b=c_{f(b)}\in\mathcal C_N,
\]
because every transformation in $T_N$ sends units to units.  Hence
$\mathcal C_N$ is a two-sided ideal.  If a nonempty ideal contains $f$, then
it contains $c_bf=c_b$ for every $b\in U_N$, so it contains
$\mathcal C_N$.  This direct proof agrees with the general minimum-rank
description for faithful finite transformation monoids
\cite[Proposition 13.1 and Lemma 13.2]{Steinberg2016}.
\end{proof}

\subsection{Reduction maps and the exact inverse limit}

For $M>N$, coefficient reduction gives
\[
 \pi_{M,N}:S_M\longrightarrow S_N,
 \qquad (\alpha,b)\longmapsto
       (\alpha\bmod3^N,b\bmod3^N).
\]

\begin{theorem}[surjectivity and exact reduction fibres]
\label{thm:finite-residue-reduction-fibres}
The maps $\pi_{M,N}$ are surjective semigroup homomorphisms.  Put $d=M-N$.
If $N\geq2$, then
\begin{equation}
\label{eq:finite-residue-general-fibres}
 |\pi_{M,N}^{-1}(s)|=
 \begin{cases}
  3^d,&s\in L_{N,1},\\
  3^{2d},&s\in L_{N,k},\quad2\leq k<N,\\
  3^{2d},&s\in\mathcal C_N.
 \end{cases}
\end{equation}
For the bottom reduction $S_M\to S_1$, each of the two fibres has size
\[
 3^{M-1}(3^{M-2}+1).
\]
The maps extend to $T_M\to T_N$, and the identity has the singleton fibre
$\{\operatorname{id}\}$.
\end{theorem}

\begin{proof}
Reduction intertwines every generator and therefore every word.  The normal
form theorem supplies a lift, proving surjectivity.  A target in $L_{N,1}$
has exactly the $3^d$ unit lifts of its intercept; its slope remains forced to
be three times that intercept.  In an intermediate stratum, the intercept
and unit slope parameter each have $3^d$ lifts, giving $3^{2d}$.

A constant target receives $3^d$ constant lifts and lifts from every source
stratum of slope valuation $k\geq N$.  Their number is
\[
 3^d\left(1+\sum_{k=N}^{M-1}\varphi_{M-k}\right)
 =3^d\left(1+\sum_{j=1}^{d}\varphi_j\right)=3^{2d}.
\]
At level one both target elements are constants.  Symmetry of the unit
intercepts modulo three splits $|S_M|$ equally, giving the stated bottom
fibre.  A nonempty word has unit intercept, so it cannot reduce to the
identity, which proves the last assertion.
\end{proof}

\begin{theorem}[exact $3$-adic inverse limit]
\label{thm:finite-residue-inverse-limit}
Under coefficientwise reduction,
\begin{equation}
\label{eq:finite-residue-inverse-limit}
 \varprojlim_N S_N
 \cong
 \{x\mapsto3b\,x+b:b\in\mathbb Z_3^\times\}
 \ \sqcup\
 \{x\mapsto\alpha x+b:
       \alpha\in9\mathbb Z_3,\ b\in\mathbb Z_3^\times\}.
\end{equation}
The second family includes $\alpha=0$.  Moreover
\[
 \varprojlim_N T_N
 =\{\operatorname{id}_{\mathbb Z_3}\}
  \sqcup\varprojlim_N S_N.
\]
For $\eta\in\mathbb Z_3^\times$ define
$G_\eta(x)=\eta(3x+1)$, and distinguish the two semigroups
\[
 \mathcal S_{\mathrm{unit}}
   =\langle G_\eta:\eta\in\mathbb Z_3^\times\rangle,
 \qquad
 \mathcal S_{\mathrm{exp}}
   =\langle G_{2^{-a}}:a\geq1\rangle.
\]
In the coefficientwise $3$-adic topology,
\[
 \overline{\mathcal S_{\mathrm{exp}}}
 =\overline{\mathcal S_{\mathrm{unit}}}
 =\varprojlim_N S_N.
\]
\end{theorem}

\begin{proof}
A compatible transformation family determines unique compatible coefficients
\[
 b_N=f_N(0),
 \qquad \alpha_N=f_N(1)-f_N(0),
\]
hence unique $b\in\mathbb Z_3^\times$ and $\alpha\in\mathbb Z_3$.
At every level $N\geq2$, the normal form says either
$\alpha_N=3b_N$ or $\alpha_N\in9R_N$.  A higher-level element of the second
kind cannot reduce to a level-$N$ element of the first kind, whose slope has
valuation exactly one.  Thus a compatible family lies entirely in the first
branch from level two onward, giving $\alpha=3b$, or entirely in the second,
giving $\alpha\in9\mathbb Z_3$.  Conversely every displayed coefficient pair
reduces to the finite normal form at every level, so it defines a compatible
family.  Composition of affine coefficient pairs is
\[
 (\alpha,b)\circ(\gamma,d)=(\alpha\gamma,\alpha d+b).
\]
In the displayed family $b,d$ are units and $3\mid\alpha$, so
$\alpha d+b$ is a unit, while $\alpha\gamma\in9\mathbb Z_3$.  Hence the
family is closed under composition, and coefficient extraction respects
composition.  This proves the asserted semigroup isomorphism.  The
identity-fibre statement in the preceding theorem prevents a mixed identity
sequence.

For arbitrary unit generators, the constructive normal-form proof works
unchanged over $\mathbb Z_3$: a word of length one gives $\alpha=3b$, and
words of each fixed length $k\geq2$ give every pair with
$\nu_3(\alpha)=k$ and unit intercept.  Closing these strata adds exactly the
zero slopes.  Thus
$\overline{\mathcal S_{\mathrm{unit}}}$ is the displayed inverse limit.

For the actual exponent generators, Lemma~\ref{lem:order-two-mod-three-power}
shows that their reductions $2^{-a}$ fill $U_N$ at every level $N$.  Given
any cylinder determined by an element of $S_N$, choose an exponent separately
for every unit letter in one representing word.  The resulting element of
$\mathcal S_{\mathrm{exp}}$ has the prescribed reduction.  Therefore
$\mathcal S_{\mathrm{exp}}$ meets every finite-level cylinder and is dense in
the same inverse limit.  This proves both closure equalities without
identifying the countable exponent semigroup with the all-unit semigroup.
\end{proof}

The slope condition cannot be weakened to
$\alpha\in3\mathbb Z_3$ with arbitrary unit intercept.  At valuation one the
exact relation is $\alpha=3b$; for example $x\mapsto3x+2$ is already excluded
modulo nine.

\subsection{Typed finite actions and the two raw obstructions}
\label{subsec:typed-finite-actions}

The specific object-erasing assignment that sends each path to its residue
transformation cannot extend linearly to an algebra representation in one
transformation algebra.  For distinct objects $x\ne y$,
\[
 e_xe_y=0,
\]
whereas erasing the object coordinate sends both empty paths to the same
identity transformation.  The product of their images is then nonzero.

There is a separate variance defect.  For a map $h:R_N\to R_N$, let
$P_hf=f\circ h$.  Then
\[
 P_hP_k=P_{k\circ h},
\]
so pullback reverses composition.  This is visible on the actual composable
path
\[
 3\xrightarrow{1}5\xrightarrow{4}1.
\]
Modulo nine,
\[
 G_1(r)=6r+5,
 \qquad G_4(r)=3r+4,
\]
but
\[
 G_4\circ G_1=c_1,
 \qquad G_1\circ G_4=c_2.
\]
Thus the two orderings are genuinely different transformations, not merely
different notation.

Let $R_0=R[u,v]$ and use the monoid algebra $R_0[T_N]$ with multiplication
induced by ordinary target-left transformation composition.  Put
\[
 \mathcal B_N=\operatorname{Mat}^{\mathrm{fin}}_{\mathcal O}(R_0[T_N]),
\]
the algebra of matrices having finite support in the infinite object index.
For the decorated prefix coordinate, put
\[
 \mathcal B_{\mathrm{pre}}
 =\operatorname{Mat}^{\mathrm{fin}}_{\mathcal O}
   \left(R_0\left[(\mathbb Z_3\rtimes_\theta P)^{\mathrm{op}}\right]\right),
\]
where the bracket denotes the finite-support monoid algebra.

\begin{proposition}[typed forward and pullback representations]
\label{prop:weighted-typed-representations}
For $p:x\to y$, let $G_p$ be its forward residue transformation.  The map
\begin{equation}
\label{eq:typed-forward-representation}
 \Pi_N([p])=w(p)E_{y,x}\otimes[G_p]
\end{equation}
is an $R_0$-algebra homomorphism
$\mathfrak A_C\to\mathcal B_N$.

On the explicitly based free module
\[
 \mathcal V_N=\bigoplus_{x\in\mathcal O}
     \bigoplus_{r\in R_N}R_0e_{x,r}
\]
it has the concrete forward action
\[
 \Pi_N([p])e_{x,r}=w(p)e_{y,G_p(r)},
\]
and is zero on every other object fibre.

For pullback, put
\[
 \mathcal F_N=\operatorname{Fun}(R_N,R_0),
 \qquad
 \mathcal K_N=
 \operatorname{Mat}^{\mathrm{fin}}_{\mathcal O}
   \bigl(\operatorname{End}_{R_0}(\mathcal F_N)\bigr).
\]
The typed pullback map
\[
 K_N([p])=w(p)E_{x,y}\otimes P_{G_p}
\]
from $\mathfrak A_C$ to $\mathcal K_N$ is an anti-homomorphism:
\[
 K_N(ab)=K_N(b)K_N(a).
\]
Finally, the decorated prefix map has the covariant typed realization
\[
 \Psi([p])
 =w(p)E_{t(p),s(p)}\otimes
   [\chi_{\mathrm{pre}}(\operatorname{word}(p))]_{\mathrm{op}}
\]
as an $R_0$-algebra homomorphism
$\Psi:\mathfrak A_C\to\mathcal B_{\mathrm{pre}}$.
\end{proposition}

\begin{proof}
If $p:x\to y$ and $q:z\to x$, then matrix-unit multiplication and forward
composition give
\[
 \Pi_N([p])\Pi_N([q])
 =w(p)w(q)E_{y,z}\otimes[G_p\circ G_q]
 =\Pi_N([p\circ q]).
\]
For noncomposable paths the matrix units multiply to zero.  Also
$\Pi_N(e_x)=E_{x,x}\otimes[\operatorname{id}]$, so the object idempotents
remain orthogonal.  This proves the algebra and concrete-module statements.

The identity $P_hP_k=P_{k\circ h}$ proves the pullback anti-homomorphism,
with the transposed matrix unit sending target functions back to source
functions.  For the last statement, the chronological word of $p\circ q$ is
the word of $q$ followed by the word of $p$.  Proposition
\ref{prop:weighted-prefix-cocycle-injection} therefore gives
\[
 \chi_{\mathrm{pre}}(\operatorname{word}(p\circ q))
 =\chi_{\mathrm{pre}}(\operatorname{word}(q))\star
  \chi_{\mathrm{pre}}(\operatorname{word}(p)),
\]
which is exactly the product in the opposite semigroup algebra after the
matrix units have enforced composability.
\end{proof}

\begin{theorem}[exact kernels of the typed representations]
\label{thm:weighted-typed-representation-kernels}
Let $R$ be a nonzero commutative ring.  The decorated prefix representation
$\Psi$ is injective.  For
$F=\sum_p f_p[p]\in\mathfrak A_C$, the finite-residue representation has the
exact kernel criterion
\begin{equation}
\label{eq:weighted-finite-representation-kernel}
 F\in\ker\Pi_N
 \quad\Longleftrightarrow\quad
 \sum_{\substack{p:x\to y\\G_p=h}}f_p w(p)=0
 \quad\text{for every }x,y\in\mathcal O\text{ and }h\in T_N.
\end{equation}
In particular, every $\Pi_N$ is noninjective.  If
$\epsilon_1:1\to1$ is the exponent-two loop, then
\begin{equation}
\label{eq:weighted-explicit-finite-kernel-element}
 uv[\epsilon_1^N]-[\epsilon_1^{N+1}]
 \in\ker\Pi_N\cap\ker K_N
\end{equation}
is nonzero.
\end{theorem}

\begin{proof}
Distinct paths have either different source-target matrix units or different
exponent words.  In the latter case the decorated map
$\chi_{\mathrm{pre}}$ separates them by
Proposition~\ref{prop:weighted-prefix-cocycle-injection}.  Hence distinct path
basis elements have distinct matrix-monoid basis coordinates under $\Psi$.
Multiplication by the monomial $w(p)$ only shifts polynomial exponents and is
injective in $R[u,v]$.  Finite linear independence proves that $\Psi$ is
injective.

For $\Pi_N$, collect the image of a finite sum by its matrix coordinate
$(y,x)$ and transformation basis element $[h]$.  Its coefficient is exactly
the sum on the right of
\eqref{eq:weighted-finite-representation-kernel}; independence of those basis
coordinates proves both directions of the criterion.

Finally $a(1)=2$, so $w(\epsilon_1)=uv$.  In $R_N$, where $4$ is invertible,
the residue transformation $H=G_{\epsilon_1}$ satisfies
\[
 H(r)-1=\frac34(r-1).
\]
Therefore $H^N(r)-1=(3/4)^N(r-1)=0$ in $R_N$, so
$H^N=H^{N+1}=c_1$.  The two terms in
\eqref{eq:weighted-explicit-finite-kernel-element} consequently have the same
image $(uv)^{N+1}E_{1,1}\otimes[c_1]$.  Their pullback images likewise agree,
because both use $P_{c_1}$ and the same matrix unit.  The original paths are
distinct basis elements and $R$ is nonzero, so their displayed difference is
nonzero.  No equality between the full kernels of $K_N$ and $\Pi_N$ is
asserted.
\end{proof}

The target $\mathcal B_N$ is finite in each residue fibre but not
finite-dimensional: it retains the infinite object set $\mathcal O$.  A
genuinely finite quotient would require a separately proved categorical
congruence.  Merely reducing object labels modulo $3^N$ identifies orthogonal
 idempotents and does not provide one.  Theorem
\ref{thm:finite-residue-semigroup-normal-form} classifies the image
transformations and their strata, not the word fibres.  The exact linear
collision criterion is instead
\eqref{eq:weighted-finite-representation-kernel}.

\subsection{What the raw programme does and does not construct}

Several later slogans in the passage fail before any Collatz-specific issue
arises.  The exact failures are useful because they determine which algebraic
structures have actually been built.

\begin{sourceissue}[four algebraic type failures]
\label{sourceissue:weighted-path-four-type-failures}
First, the natural reduction maps
\[
 \mathbb Z/3^{N+1}\mathbb Z\longrightarrow\mathbb Z/3^N\mathbb Z
\]
form an inverse system, not the direct system printed in an early draft.  Its
limit is $\mathbb Z_3=\varprojlim_N\mathbb Z/3^N\mathbb Z$; a residue class
modulo $3^m$ has no canonical lift modulo $3^{m+n}$.

Second, $\mathbb R_{>0}$ under its usual addition and multiplication is not a
semiring because it has no additive zero.  The path weights may instead take
values in the multiplicative monoid $\mathbb R_{>0}$ or in the semiring
$\mathbb R_{\geq0}$.

Third, ordinary deconcatenation is a coalgebra coproduct on words but is not
multiplicative for ordinary concatenation.  For two letters $a,b$,
\[
 \Delta(a)\Delta(b)
 =1\otimes ab+a\otimes b+b\otimes a+ab\otimes1,
\]
whereas
\[
 \Delta(ab)=1\otimes ab+a\otimes b+ab\otimes1.
\]
The uncancelled term $b\otimes a$ proves that the printed data do not define
a bialgebra.  An unshuffle coproduct, a braided tensor product, or another
explicit compatibility law would be a different construction.

Fourth, the two-element set of local affine formulas
$\{X/2,(3X+1)/2\}$ is not closed under affine composition: already
$(X/2)\circ(X/2)=X/4$.  It is therefore not an algebra acting on the path
algebra, and no crossed product with it is defined by the passage.
\end{sourceissue}

At the proved scope, the programme supplies:
\begin{enumerate}[label=\textup{(\arabic*)}]
\item an exact equivalence between exponent paths and legal local
odd-to-odd blocks;
\item a genuine two-variable path weight and two explicitly compared
$\mathbb F_2^2$ degrees;
\item endpoint formal series, a coefficientwise path completion, and exact
resolvent identities;
\item a predecessor transfer operator with precise algebraic, formal, and
analytic domains;
\item an injective decorated prefix map $\chi_{\mathrm{pre}}$ and its exact
reversal relation to the forward affine/Tao offset;
\item finite residue transformation semigroups with complete normal forms,
cardinalities, ranks, minimum ideals, reduction fibres, and inverse limit;
\item typed forward representations and pullback anti-representations that
retain object idempotents and composition order.
\end{enumerate}

\begin{nonclaim}
The directly read passage does not construct a Collatz path groupoid of local
homeomorphisms, a groupoid $C^*$-algebra, a Leavitt or Cuntz--Krieger algebra,
a quotient category $\operatorname{QGr}$, a Hopf or bialgebra structure, a
crossed coefficient product, an operator-algebra completion, a BCM system, an
adelic system, a KMS state, or a uniqueness or canonicity theorem for any such
object.  Names and analogies appearing after the finite constructions do not
supply the missing objects or maps.  The rejected BCM/adelic material in the
later local AI-assisted PDF is not imported here.  None of the theorems in
this section proves the Collatz conjecture, and no priority claim is made for
the independently reconstructed results.
\end{nonclaim}

The finite certificate
\path{certificates/chatnotes_weighted_path_finite_actions.py} pins the raw
passage, all three human references, both local AI-assisted lineage witnesses,
Tao's v7 source, and the existing literature bridge.  It verifies bounded
instances of the block inverse, grading fibres, cocycle and reversal formulas,
the noninjective undecorated $\rho$ example, source congruences, predecessor
fibres, endpoint coefficients, finite path-cut counts, the exact generated
residue semigroups through $N=6$, their ranks and adjacent reduction fibres,
finite projections of the inverse-limit classification, typed path products,
and the explicit fixed-point-loop kernel relation.  It does not independently
prove the infinite, analytic, faithfulness, or exact-kernel theorems printed
above.
